% Generated draft: requires mathematical and source review.
\documentclass[11pt]{article}
\usepackage[margin=25mm]{geometry}
\usepackage{fontspec}
\setmainfont{Latin Modern Roman}
\usepackage{amsmath,amssymb,mathtools,unicode-math}
\setmathfont{Latin Modern Math}
\usepackage{enumitem,needspace,xurl,hyperref}
\setlength{\parindent}{0pt}
\setlength{\parskip}{5pt}
\setlength{\emergencystretch}{3em}
\newcommand{\N}{\mathbb{N}}
\newcommand{\Z}{\mathbb{Z}}
\newcommand{\Q}{\mathbb{Q}}
\newcommand{\R}{\mathbb{R}}
\newcommand{\C}{\mathbb{C}}
\newcommand{\F}{\mathbb{F}}
\newcommand{\A}{\mathbb{A}}
\newcommand{\PP}{\mathbb{P}}
\newcommand{\E}{\mathbb{E}}
\newcommand{\eps}{\varepsilon}
\newcommand{\dd}{\,\mathrm{d}}
\begin{document}
\section*{UnsolvedMath ID: 30005433}
\textbf{Source problem code:} catalogue entry
\subsection{Definitions and mathematical statement}
A (left) skew brace is a set \(B\) with group operations \(+\) and \(\circ\), having the same identity \(0\), such that
\[
a\circ(b+c)=(a\circ b)-a+(a\circ c)\qquad(a,b,c\in B).
\]
Neither group is assumed abelian. Put
\[
\lambda_a(b)=-a+(a\circ b),\qquad a*b=-a+(a\circ b)-b.
\]
An ideal of \(B\) means a subset normal in both groups and invariant under every \(\lambda_a\), \(a\in B\). For \(X\subseteq B\), let \(\langle X\rangle_{\mathrm{ideal},B}\) denote the least ideal of \(B\) containing \(X\). In the formulations below, ``finitely generated ideal'' is interpreted as \(\langle X\rangle_{\mathrm{ideal},B}\) for some finite \(X\). Generation as a skew brace means generation under both group operations and their inverses; it is a distinct requirement.

For \(x\in B\), define
\[
\begin{aligned}
\operatorname{Fix}^{r}(x)&=\{b\in B:x*b=0\},&
\operatorname{Fix}^{l}(x)&=\{b\in B:b*x=0\},\\
C_x^+&=\{b\in B:b+x=x+b\},&
C_x^\circ&=\{b\in B:b\circ x=x\circ b\}.
\end{aligned}
\]
The brace has property \((S)\) if, for every \(x\in B\),
\[
[(B,+):\operatorname{Fix}^{r}(x)\cap C_x^+]<\infty,
\qquad
[(B,\circ):\operatorname{Fix}^{l}(x)\cap C_x^\circ]<\infty.
\]
Under the ambient-ideal interpretation above, call an ideal \(I\) good if every ideal \(J\) of \(B\) contained in \(I\) that is finitely generated as an ideal of \(B\) is finitely generated as a skew brace. In particular, goodness of \(B\) asks whether
\[
\forall X\subseteq B\text{ finite},\quad
\langle X\rangle_{\mathrm{ideal},B}
\text{ is finitely generated as a skew brace}.
\]

The supplied cleaned record asks the following two questions for every skew brace \(B\) with property \((S)\):
\begin{enumerate}
\item Is \(\langle X\rangle_{\mathrm{ideal},B}\) finite for every finite subset \(X\subseteq B\)?
\item Is \(B\) good?
\end{enumerate}
These are retained as separate source formulations; their equivalence is not asserted in this draft.

\textit{Unresolved source qualifications.} The original passage also asks whether every \emph{periodic} skew brace with property \((S)\) is good, and attributes an equivalence with the unrestricted goodness question to Theorems 3.14, 3.25 and 3.27 of its reference [1]. The supplied record neither identifies [1] nor reproduces those theorems. It also does not specify which underlying group is required to be periodic, where a periodic group is one in which every element has finite order. The first finiteness question in the original passage refers to a previously specified \(B\), whose complete preceding hypotheses are not supplied. Whether periodicity or another hypothesis belongs to that finiteness question therefore requires review. The ambient interpretation of ideal generation used above must likewise be checked against the source, particularly for goodness of a proper ideal \(I\).
\subsection{Short English statement}
The problem asks whether ideals generated by finitely many elements in a skew brace with property \((S)\) must be finite, and whether such ideals must be finitely generated as skew braces. The original passage also includes a periodicity restriction and a claimed equivalence whose hypotheses and cited justification are not fully supplied. These formulations require human mathematical review.
\subsection{Sources}
\begin{itemize}
\item Local registry: \texttt{data/unsolvedmath/problems.json}, permanent UnsolvedMath ID \(30005433\), source problem code \texttt{catalogue entry}, title: ``Finiteness and Goodness for Property-S Skew Braces.'' The supplied record gives proposed year 2023 and no individual proposer.
\item Source report: \emph{Mini-Workshop: Skew Braces and the Yang--Baxter Equation} (2023), Oberwolfach Reports, pp. 537--563, DOI 10.4171/owr/2023/9. Supplied source URL and literature-source link: \url{https://doi.org/10.4171/owr/2023/9}.
\item Supplied literature source: Colazzo--Ferrara--Trombetti, \emph{Publ. Mat.} 69 (2025), \url{https://mat.uab.es/pubmat/fitxers/download/FileType%3Apdf/FolderName%3Av69%281%29/FileName%3A6912508.pdf}. No article title is supplied.
\item Literature assessment dated \textbf{2026-08-22}: the supplied curation labels the problem open and states that the finite-generation question for ideals, described there as equivalent to goodness for every skew brace with property \((S)\), remains explicitly open in the 2025 treatment. This is supplied dated curation, not an independently verified literature assessment.
\item The registry labels the statement \texttt{corrected\_verified} and reports a direct check of the official report. Those are supplied provenance claims; no source visit or independent verification is claimed here. The original passage cites Theorems 3.14, 3.25 and 3.27 of reference [1], but the bibliographic identity of [1] is absent from the supplied record.
\end{itemize}
\end{document}
